% ============================================================
% Section 4: Experiments and Results
% ============================================================

\section{Experiments and Results}
\label{sec:experiments}

\subsection{Experimental Setup}
\label{subsec:setup}

All experiments are conducted on a single workstation running Windows~11,
equipped with an AMD Ryzen~5 3400G (3.70\,GHz), 32\,GB DDR4-3200 RAM, and an NVIDIA
RTX~5060\,Ti (16\,GB VRAM), using PyTorch 2.10.0 with CUDA 13.0 and
Optuna~3.x. The
full NCV pipeline required approximately 59~hours; images are cached in GPU
memory per outer fold, reducing per-trial training time from $\sim$12 to
$\sim$3.5 minutes.

\subsection{Nested Cross-Validation Results}
\label{subsec:ncv_results}

Table~\ref{tab:nested_cv} reports per-fold test metrics. PR-AUC is
consistently high across all folds ($0.977$--$0.999$, mean $0.9921 \pm
0.0058$). Dice shows slightly more fold-to-fold variability ($0.800$--$0.909$,
mean $0.856 \pm 0.030$), reflecting the difficulty of pixel-level localization.
The Pareto-selected $\lambda_{\text{mask}}$ clusters between 2.3 and 2.7 in
eight of ten folds; the two outlier folds (F5, F9) correspond to the two lowest
Dice scores, confirming that $\lambda_{\text{mask}}$ is a decisive driver of
localization quality.

\begin{table}[!t]
  \centering
  \caption{Per-fold test results from 10-fold NCV.
           $\overline{w}_{\text{mask}}$: Pareto-selected loss weight.}
  \label{tab:nested_cv}
  \begin{tabular}{cccc}
    \toprule
    Fold & PR-AUC & Dice & $\overline{w}_{\text{mask}}$ \\
    \midrule
    F1  & 0.9939 & 0.861 & 2.46 \\
    F2  & 0.9769 & 0.871 & 2.70 \\
    F3  & 0.9993 & 0.909 & 2.47 \\
    F4  & 0.9915 & 0.858 & 2.43 \\
    F5  & 0.9948 & 0.827 & 1.41 \\
    F6  & 0.9936 & 0.874 & 2.89 \\
    F7  & 0.9940 & 0.861 & 2.67 \\
    F8  & 0.9967 & 0.887 & 2.29 \\
    F9  & 0.9896 & 0.807 & 1.95 \\
    F10 & 0.9909 & 0.800 & 2.65 \\
    \midrule
    \textbf{Mean} & \textbf{0.9921} & \textbf{0.856} & \textbf{2.39} \\
    \textbf{Std}  & \textbf{0.0058} & \textbf{0.030} & \textbf{0.43} \\
    \bottomrule
  \end{tabular}
\end{table}

\subsection{Ablation Study}
\label{subsec:ablation}

To isolate each design decision, we evaluate four variants on the blind test
set using 30 independent seeds: \textit{Multitask} (ours, joint detection +
localization with Pareto-optimized $\lambda_{\text{mask}}$);
\texttt{cls\_only} (classification head only); \texttt{seg\_only}
(segmentation head only, detection by thresholding mean mask probability);
and \texttt{unweighted\_losses} (full multitask, $\lambda_{\text{cls}} =
\lambda_{\text{mask}} = 0.5$, bypassing HPO).
Table~\ref{tab:ablation} and Fig.~\ref{fig:ablation} summarize results.

\begin{figure*}[!t]
  \centering
  \includegraphics[width=\textwidth]{fig3_ablation}
  \caption{Ablation study: violin plots of PR-AUC (left) and Dice (right) for
    four variants over 30 seeds. Significance vs.\ Multitask: *** $p < 0.001$;
    ns $p \geq 0.05$ (Wilcoxon + Holm).}
  \label{fig:ablation}
\end{figure*}

\begin{table*}[!t]
  \centering
  \caption{Ablation study (blind test, 30 seeds).
           $p$: Wilcoxon + Holm; $d$: Cohen's~$d$.
           *** $p < 0.001$; ns $p \geq 0.05$.}
  \label{tab:ablation}
  \begin{tabular}{lcccccc}
    \toprule
    & \multicolumn{2}{c}{PR-AUC} & \multicolumn{4}{c}{Dice} \\
    \cmidrule(r){2-3} \cmidrule(l){4-7}
    Variant & Mean $\pm$ Std & vs.\ Multitask & Mean $\pm$ Std & vs.\ Multitask & $p_{\text{holm}}$ & $d$ \\
    \midrule
    \textbf{Multitask} (ours) & $\mathbf{0.9967 \pm 0.0021}$ & — & $\mathbf{0.875 \pm 0.018}$ & — & — & — \\
    cls\_only          & $0.9503 \pm 0.0744$ & *** & $0.002 \pm 0.008$ & *** & $1.3 \times 10^{-8}$ & large \\
    seg\_only          & $0.7277 \pm 0.0528$ & *** & $0.867 \pm 0.023$ & ns  & $0.288$              & small \\
    unweighted\_losses & $0.9958 \pm 0.0016$ & ns  & $0.857 \pm 0.016$ & *** & $1.3 \times 10^{-4}$ & 1.01 \\
    \bottomrule
  \end{tabular}
\end{table*}

Three findings stand out. First, joint multitask training is essential:
\texttt{cls\_only} produces near-zero Dice ($0.002 \pm 0.008$). Second, a
dedicated classification head is necessary: \texttt{seg\_only} achieves
competitive Dice ($0.867$) but PR-AUC drops to $0.728 \pm 0.053$, confirming
that the segmentation logit alone is an unreliable detection signal. Third,
Pareto-optimized weights outperform equal weights specifically on localization:
\texttt{unweighted\_losses} matches Multitask on PR-AUC ($p = 0.083$) but
shows significantly lower Dice ($p = 1.3 \times 10^{-4}$, $d = 1.01$, large
effect), validating the HPO contribution.

\subsection{Comparison with Related Systems}
\label{subsec:challenge}

Table~\ref{tab:challenge} compares DocVerify against the top participants of
the DeepID 2025 Challenge on Track~1 (detection F1@0.5) and Track~2
(per-image localization F1), computed on our 15\% holdout using 30 seeds.

\begin{table*}[!t]
  \centering
  \caption{Comparison with related systems. $\dagger$: official challenge test
    set; DocVerify evaluated on internal holdout (see text).}
  \label{tab:challenge}
  \begin{tabular}{lccc}
    \toprule
    System & F1 det (thr=0.5) & F1 loc & Training Data \\
    \midrule
    Sunlight~\citep{remtkd2025} (1st)$^\dagger$      & 0.991 & 0.784 & 60K+ external images \\
    Incode (2nd)$^\dagger$                            & 0.868 & —     & 100K proprietary \\
    AG/EdgeDoc~\citep{edgedoc2025} (3rd)$^\dagger$   & 0.958 & 0.686 & FantasyID only \\
    UAM-Biometrics (4th)$^\dagger$                   & 0.712 & 0.620 & FantasyID only \\
    TruFor~\citep{trufor2023} baseline$^\dagger$     & 0.807 & 0.590 & 876K pre-training \\
    \midrule
    \textbf{DocVerify (ours)} & $\mathbf{0.969 \pm 0.014}$ & $\mathbf{0.807 \pm 0.096}$ & FantasyID only \\
    \bottomrule
  \end{tabular}
\end{table*}

DocVerify achieves F1@0.5 $= 0.969 \pm 0.014$ on Track~1, surpassing all
FantasyID-only systems including AG/EdgeDoc ($+1.1$~pp) and outperforming
Incode despite its use of 100K proprietary images. On Track~2, DocVerify
achieves F1$_{\text{loc}} = 0.807 \pm 0.096$, exceeding all four challenge
finishers including the overall winner Sunlight ($0.784$). The high Track~2
variance ($\pm 0.096$) is driven by seeds where bonafide misclassifications
generate false-positive mask pixels; F1 on attack images is stable
($0.881 \pm 0.015$).

The comparison should be interpreted with caution. DocVerify was developed
independently after the DeepID 2025 Challenge submission deadline had already
closed; consequently, official test-set evaluation was not possible and we
report results on our internal 15\% holdout instead. Although both sets are
drawn from the same FantasyID distribution, they are not identical, making the
comparison indicative rather than strictly controlled.

\subsection{Qualitative Analysis}
\label{subsec:qualitative}

Figure~\ref{fig:qualitative} shows a representative output for a forged
document with five altered regions. The predicted mask successfully localizes
all five tampered fields — photograph, given name, date of birth, document
number, and expiry date — with tight spatial alignment to the ground-truth
alterations.

\begin{figure*}[!t]
  \centering
  \begin{tabular}{@{}ccc@{}}
    \small\textbf{Input document} &
    \small\textbf{Predicted mask} &
    \small\textbf{Overlay} \\[2pt]
    \includegraphics[width=0.318\textwidth]{fig_qual_attack_input} &
    \includegraphics[width=0.318\textwidth]{fig_qual_attack_mask}   &
    \includegraphics[width=0.318\textwidth]{fig_qual_attack_overlay} \\[1pt]
    \multicolumn{3}{c}{\small Forged document — five altered regions localized} \\
  \end{tabular}
  \caption{Qualitative result on a forged FantasyID sample. \textbf{Left:}
  input document. \textbf{Centre:} binary mask predicted by DocVerify.
  \textbf{Right:} predicted mask overlaid in red on the input. All five
  tampered fields are correctly localized.}
  \label{fig:qualitative}
\end{figure*}

